% ch5_c 第5章 §5.3尾–§5.4.3 (书页 213–225 / p228–p240)
%--C-TAIL-- 两个电荷由于来自其他电子的屏蔽，并不感受到长程相互作用。

%JOIN%
与具有相同动量的粒子--空穴激发的能量范围一致（见图 4.11）。事实上，该本征模对应的正是图 5.5 中的粒子--空穴激发。

为求 $f(\theta,\theta')$ 非零时集体模式的能量，我们注意到式 (5.3.14) 可以改写为
\begin{equation*}
\mathrm{i}q^{-1}\partial_{t}\tilde{\rho}(\theta,\pmb{q},t)=KM\tilde{\rho}(\theta,\pmb{q},t) \tag{5.3.15}
\end{equation*}
其中两个实对称算符 $K$ 与 $M$ 由下式给出：
\begin{align*}
K(\theta,\theta') &= \cos(\theta_{\pmb{q}}-\theta)\,\delta(\theta-\theta')\\
M(\theta,\theta') &= v_{F}^{*}\,\delta(\theta-\theta')+\frac{k_{F}}{(2\pi)^{2}}f(\theta,\theta') \tag{5.3.16}
\end{align*}
由费米液体的能量（见式 (5.3.12)）可知，费米液体的稳定性要求 $M$ 是正定的。于是我们可以把 $M$ 写成 $M=\tilde{M}\tilde{M}^{\top}$。令 $u=\tilde{M}^{\top}\tilde{\rho}$，则式 (5.3.15) 变为 $\mathrm{i}\partial_{t}u=q\tilde{M}^{\top}K\tilde{M}u$。具有固定动量 $\pmb{q}$ 的集体激发，其能量由实对称矩阵 $q\tilde{M}^{\top}K\tilde{M}$ 的本征值给出。

转动对称性要求 $f(\theta,\theta')=f(\theta-\theta')$。由于能谱不依赖于 $\pmb{q}$ 的方向，我们可以取 $\pmb{q}$ 沿 $x$ 方向，即 $\pmb{q}=q\pmb{x}$、$\theta_{\pmb{q}}=0$。引入（见图 5.4）
\begin{equation*}
\tilde{\rho}_{l}(\pmb{q},t)=\int \mathrm{d}\theta\,\frac{\mathrm{e}^{-\mathrm{i}l\theta}}{\sqrt{2\pi}}\,\tilde{\rho}(\theta,\pmb{q},t),
\end{equation*}
式 (5.3.15) 便可写成
\begin{align*}
\mathrm{i}\partial_{t}\tilde{\rho}_{l}(q\pmb{x},t) &= qK_{lm}M_{ml'}\tilde{\rho}_{l'}(q\pmb{x},t), & f(\theta-\theta') &= \sum_{l}f_{l}\,\mathrm{e}^{\mathrm{i}l(\theta-\theta')},\\
K_{ll'} &= \frac{1}{2}\left(\delta_{l,l'+1}+\delta_{l,l'-1}\right), & M_{ll'} &= \left(v_{F}^{*}+\frac{k_{F}f_{l}}{2\pi}\right)\delta_{ll'}.
\end{align*}
由 $M$ 可见，若某一个 $f_{l}$ 小于 $-2\pi v_{F}^{*}/k_{F}$，费米液体就是不稳定的。我们还看到 $\tilde{M}_{ll'}\allowbreak=\sqrt{v_{F}^{*}+\frac{k_{F}f_{l}}{2\pi}}\,\delta_{ll'}$，以及
\begin{equation*}
q(\tilde{M}^{\top}K\tilde{M})_{ll'}\equiv\tilde{H}_{ll'}=\frac{q}{2}\left(\delta_{l,l'+1}+\delta_{l,l'-1}\right)\sqrt{v_{F}^{*}+\frac{k_{F}f_{l}}{2\pi}}\sqrt{v_{F}^{*}+\frac{k_{F}f_{l'}}{2\pi}}
\end{equation*}
$\tilde{H}$ 的本征值给出集体模式的频率。

这里 $\tilde{H}$ 描述一个在一维格子上跳跃的粒子。位点 $l$ 与位点 $l+1$ 之间的最近邻跳跃幅为
$\frac{q}{2}\sqrt{\frac{v_{F}^{*}}{2\pi}+\frac{k_{F}f_{l}}{(2\pi)^{2}}}\sqrt{\frac{v_{F}^{*}}{2\pi}+\frac{k_{F}f_{l+1}}{(2\pi)^{2}}}$。

%FIG{5.6}: {(a) $f_{l}=0$ 时 $\tilde{H}$ 的谱，(b) 相应的粒子--空穴谱。(c) $f_{l}>0$ 时 $\tilde{H}$ 的谱，(d) 相应的粒子--空穴谱。}

当 $f_{l}=0$ 时，$\tilde{H}$ 具有从 $-v_{F}^{*}q$ 到 $v_{F}^{*}q$ 的连续谱（见图 5.6(a)），它重现了自由费米子系统的粒子--空穴连续谱（见图 5.6(b)）。然而，若 $f_{l}>f_{\infty}$，则由 $\tilde{H}_{ll'}$ 所描述的粒子会被吸引到小 $l$ 的区域。于是，除连续谱之外，$\tilde{H}$ 的谱中还有对应于吸引区内束缚态的孤立本征值（见图 5.6(c)）。相应的粒子--空穴谱则包含轮廓分明的集体模式（见图 5.6(d)）。\footnote{玻色型模式必须具有正能量，负能量意味着不稳定。然而，$\tilde{H}$ 既有正本征值也有负本征值，分别对应正、负频率。如果 $\tilde{H}$ 的正本征值对应玻色集体模式的正能量，那么 $\tilde{H}$ 的负本征值是否也对应集体模式的负能量？这里我们想指出，到目前为止我们一直是把式 (5.3.13) 当作经典拉格朗日量来处理的，也只讨论了由此得到的经典运动方程。量子化之后，运动方程变成关于算符 $\tilde{\rho}$ 的方程。事实表明，正频率的模式对应玻色模式的产生算符，而负频率的模式对应玻色模式的湮灭算符。式 (5.3.13) 的一维版本的量子化将在 7.4.5 节讨论，见式 (7.4.35) 与式 (7.4.36)。}

在第二个应用中，我们想用线性化的流体动力学运动方程 (5.3.10) 来计算费米液体的直流电导。当存在均匀的静态电场 $\pmb{E}$ 时，力由 $\pmb{F}=e\pmb{E}$ 给出。力 $\pmb{F}$ 将诱导出一个均匀且静态的 $\tilde{\rho}$，即 $\partial_{t}\tilde{\rho}=\partial_{\pmb{x}}\tilde{\rho}=0$。式 (5.3.10) 变为
\begin{equation*}
\tilde{\rho}=e\tau\,\frac{k_{F}}{(2\pi)^{2}}\,\hat{\pmb{k}}\cdot\pmb{E}
\end{equation*}
由式 (5.3.6) 可见，上述 $\tilde{\rho}$ 诱导出电流：
\begin{align*}
e\pmb{J} &= e\int \mathrm{d}\theta\,\tilde{\rho}\,\tilde{\pmb{v}}\\
\tilde{\pmb{v}}(\pmb{k}) &= \pmb{v}^{*}(\pmb{k})+\int \frac{\mathrm{d}^{2}\pmb{k}'}{(2\pi)^{2}}\,f(\pmb{k},\pmb{k}')\,\pmb{v}^{*}(\pmb{k}')\,\delta(\xi_{\pmb{k}'}^{*})
\end{align*}
由于转动对称性，$\tilde{\pmb{v}}=|\tilde{\pmb{v}}|\hat{\pmb{k}}$。我们得到
\begin{equation*}
e\pmb{J}=e^{2}\tau\,\frac{|\tilde{\pmb{v}}|k_{F}}{(2\pi)^{2}}\int \mathrm{d}\theta\,\hat{\pmb{k}}(\hat{\pmb{k}}\cdot\pmb{E})=\frac{e^{2}\tau|\tilde{\pmb{v}}|\rho}{k_{F}}\,\pmb{E}
\end{equation*}
电导为 $\sigma=\frac{e^{2}\tau|\tilde{\pmb{v}}|\rho}{k_{F}}$。若费米液体函数 $f(\pmb{k},\pmb{k}')$ 为零，上式就退化为 Drude 结果 $\sigma=\frac{e^{2}\tau\rho}{m^{*}}$。

\subsection*{问题 5.3.5}

当存在均匀磁场 $\pmb{B}=B\pmb{z}$ 时，二维约化玻尔兹曼方程 (5.3.9) 中的力 $\pmb{F}$ 具有 $\pmb{F}(\pmb{k},\pmb{x})=\frac{e}{c}\pmb{v}^{*}(\pmb{k})\times\pmb{B}$ 的形式。
\begin{enumerate}
\item[(a)] 证明：线性化之后，磁场中的约化玻尔兹曼方程在取 $\tau^{-1}=0$ 时具有如下形式（Kim et al., 1995）：
\begin{equation*}
\left(\mathrm{i}\partial_{t}+\mathrm{i}\omega_{c}\partial_{\theta}-v_{F}^{*}q\cos(\theta-\theta_{\pmb{q}})\right)\tilde{\rho}-\frac{qk_{F}\cos(\theta_{\pmb{q}}-\theta)}{(2\pi)^{2}}\int \mathrm{d}\theta'\,f(\theta-\theta')\,\tilde{\rho}(\theta',\pmb{q},t)=0
\end{equation*}
试求出 $\omega_{c}$。
\item[(b)] 证明：动量为 $\pmb{q}$ 的集体模式的能谱由下式的本征值决定：
\begin{equation*}
\tilde{H}_{ll'}=\omega_{c}l\,\delta_{ll'}+\frac{q}{2}\left(\delta_{l,l'+1}+\delta_{l,l'-1}\right)\sqrt{v_{F}^{*}+\frac{k_{F}f_{l}}{2\pi}}\sqrt{v_{F}^{*}+\frac{k_{F}f_{l'}}{2\pi}}
\end{equation*}
证明 $\tilde{H}$ 的谱关于零对称。（提示：$f_{l}=f_{-l}$ 为实数。）
\item[(c)] 证明：当 $f_{l}=0$ 时，$\tilde{H}$ 的本征值由 $\omega_{c}\times$ 整数给出。（提示：考虑 $\tilde{\rho}=\exp\bigl(\mathrm{i}l\theta-\mathrm{i}\frac{v_{F}^{*}q}{\omega_{c}}\sin(\theta-\theta_{\pmb{q}})\bigr)$。）这样我们便能用流体动力学理论恢复自由费米子的结果。
\end{enumerate}

\subsection*{问题 5.3.6}

\textbf{二维费米液体的交流输运}

\begin{enumerate}
\item[(a)] 设式 (5.3.9) 中的力 $\pmb{F}$ 具有 $\pmb{F}=e\pmb{E}\,\mathrm{e}^{\mathrm{i}\omega t-\pmb{q}\cdot\pmb{x}}$ 的形式。在 $f(\theta-\theta')=0$ 的假设下求 $\tilde{\rho}$。
\item[(b)] 准确到 $f(\theta-\theta')$ 的一阶，求 $\tilde{\rho}$。
\item[(c)] 求感应电流 $e\pmb{J}=\sigma(\omega,\pmb{q})\pmb{E}$，并准确到 $f(\theta-\theta')$ 的一阶求光学电导 $\sigma(\omega,\pmb{q})$。这里 $\pmb{J}$ 是式 (5.3.6) 给出的准粒子电流。
\end{enumerate}

\subsection{费米液体理论的本质}

人们通常认为，费米液体理论的精髓在于轮廓分明的准粒子这一概念及其实际存在。在 5.3.1 节中我们采纳的正是这一观点。轮廓分明的准粒子的存在意味着，不仅准粒子的总数守恒，而且在每个给定动量方向上的准粒子数目也守恒。于是，费米液体具有无穷多个守恒量（Luther, 1979; Haldane, 1992）。这使我们能够为费米液体建立一套流体动力学理论，其中包含许多玻色型模式，每个守恒量对应一个模式。因此，无穷多守恒量的存在也是费米液体理论的精髓之一。在 5.3.3 节中，我们正是基于这一观点建立了费米液体理论。

我们想指出，这两种观点并不等价，而且第二种观点更为普遍。轮廓分明的准粒子的存在意味着每个动量上的准粒子数目守恒；而在第二种观点中，我们只假设每个动量方向上的准粒子守恒。在一维情形，5.3.3 节的流体动力学理论实际上描述的是一维 Tomonaga--Luttinger 液体。\footnote{关于 Tomonaga--Luttinger 液体的讨论见 7.4.4 节。}在更高维数下，流体动力学理论描述的则是费米液体。

\section{微扰论与费米液体理论的有效性}

本节我们将检验费米液体理论的正确性。我们将发展一套微扰论来计算相互作用系统的电子格林函数，并检查 $\operatorname{Im}G(\pmb{k},\omega)$ 中是否含有 $\delta$ 函数。$\delta$ 函数的存在标志着轮廓分明的 Landau 准粒子的存在，也标志着 Landau 费米液体理论的适用。

\subsection{费米子的路径积分与微扰论}

\begin{keypoints}
\item 费米子的路径积分是一种记账形式体系，它使我们能够把费米子系统的关联函数形式地表达出来。与玻色子的路径积分不同，费米子路径积分的物理含义并不清楚。
\item 费米子的路径积分使我们能够为相互作用费米子系统系统地发展微扰论。
\item 微扰论的结构由费曼图与费曼规则所刻画。
\end{keypoints}

为建立费米子的路径积分，我们首先注意到，费米子算符的编时关联满足 $\langle T(c(t)c^{\dagger}(t'))\rangle=-\langle T(c^{\dagger}(t')c(t))\rangle$。于是，我们将用（复的）格拉斯曼数（Grassmann number）$\xi_{\alpha}$ 来表示费米子算符。读者可能会问，``用格拉斯曼数表示费米子算符''是什么意思？嗯，我不得不承认我不知道。我不知道格拉斯曼数的物理含义。我将把下面的费米子路径积分当作某种记账工具，它使我们能够把各种费米子关联函数形式地装进一个紧凑的公式里。格拉斯曼数是反对易数，满足
\begin{equation*}
\xi_{\alpha}\xi_{\beta}=-\xi_{\beta}\xi_{\alpha},\qquad \xi_{\alpha}\xi_{\beta}^{*}=-\xi_{\beta}^{*}\xi_{\alpha},\qquad \xi_{\alpha}^{2}=0,
\end{equation*}
以及
\begin{equation*}
\left(\xi_{\alpha_{1}}\dots\xi_{\alpha_{n}}\right)^{*}=\xi_{\alpha_{n}}^{*}\dots\xi_{\alpha_{1}}^{*}
\end{equation*}
格拉斯曼数的函数由如下展开定义：
\begin{equation*}
f(\xi)=f_{0}+f_{1}\xi,\qquad A(\xi,\xi^{*})=a_{0}+a_{1}\xi+\bar{a}_{1}\xi^{*}+a_{12}\xi^{*}\xi
\end{equation*}
导数的定义与普通导数相同，唯一的区别是：为了使导数算符 $\frac{\partial}{\partial\xi}$ 能作用到 $\xi$ 上，必须通过反对易把变量 $\xi$ 一路移到与 $\frac{\partial}{\partial\xi}$ 相邻的位置。例如，$\frac{\partial}{\partial\xi}(\xi^{*}\xi)=\frac{\partial}{\partial\xi}(-\xi\xi^{*})=-\xi^{*}$。在这些定义下，我们有
\begin{align*}
\frac{\partial}{\partial\xi}A(\xi,\xi^{*}) &= a_{1}-a_{12}\xi^{*}\\
\frac{\partial}{\partial\xi^{*}}A(\xi,\xi^{*}) &= \bar{a}_{1}+a_{12}\xi\\
\frac{\partial}{\partial\xi^{*}}\frac{\partial}{\partial\xi}A(\xi,\xi^{*}) &=-a_{12}=-\frac{\partial}{\partial\xi}\frac{\partial}{\partial\xi^{*}}A(\xi,\xi^{*})
\end{align*}
我们看到
\begin{equation*}
\frac{\partial}{\partial\xi^{*}}\frac{\partial}{\partial\xi}=-\frac{\partial}{\partial\xi}\frac{\partial}{\partial\xi^{*}}
\end{equation*}
格拉斯曼积分作为一种线性映射，被定义为格拉斯曼导数：
\begin{equation*}
\int \mathrm{d}\xi=\frac{\partial}{\partial\xi}
\end{equation*}
我们有
\begin{equation*}
\int \mathrm{d}\xi\,1=0,\qquad \int \mathrm{d}\xi\,\xi=1
\end{equation*}
格拉斯曼积分具有标准积分的如下性质：
\begin{equation*}
\int \mathrm{d}\xi\,\frac{\partial}{\partial\xi}f(\xi)=0,\qquad \int \mathrm{d}\xi\,f(\xi+\eta)=\int \mathrm{d}\xi\,f(\xi)
\end{equation*}
从行列式的定义
\begin{equation*}
\mathrm{Det}(H_{0})=\sum_{P}(-)^{P}(H_{0})_{1P_{1}}\dots(H_{0})_{nP_{n}}
\end{equation*}
（其中 $P$ 为置换 $(1,\dots,n)\to(P_{1},\dots,P_{n})$）出发，我们得到格拉斯曼高斯积分：
\begin{equation*}
\int \prod_{i=1}^{n}\mathrm{d}\eta_{i}^{*}\,\mathrm{d}\eta_{i}\,\mathrm{e}^{-\eta_{i}^{*}(H_{0})_{ij}\eta_{j}}=\mathrm{Det}(H_{0})
\end{equation*}
以及
\begin{equation*}
\int \prod_{i=1}^{n}\mathrm{d}\eta_{i}^{*}\,\mathrm{d}\eta_{i}\,\mathrm{e}^{-\eta_{i}^{*}(H_{0})_{ij}\eta_{j}+\zeta_{i}^{*}\eta_{i}+\zeta_{i}\eta_{i}^{*}}=\mathrm{Det}(H_{0})\,\mathrm{e}^{\zeta_{i}^{*}(H_{0}^{-1})_{ij}\zeta_{j}}
\end{equation*}
注意，与玻色型高斯积分相比，行列式出现在分子而不是分母中。

费米型高斯积分使我们能够得到（由路径积分定义的）费米子关联的 Wick 定理：
\begin{align*}
&\left\langle \psi_{1}\dots\psi_{n}\psi_{n}^{*}\dots\psi_{1}^{*}\right\rangle\\
&=\frac{\int \mathcal{D}\psi^{*}\mathcal{D}\psi\,\psi_{1}\dots\psi_{n}\psi_{n}^{*}\dots\psi_{1}^{*}\,\mathrm{e}^{-\sum_{ij}\psi_{i}^{*}(H_{0})_{ij}\psi_{j}}}{\int \mathcal{D}\psi^{*}\mathcal{D}\psi\,\mathrm{e}^{-\sum_{ij}\psi_{i}^{*}(H_{0})_{ij}\psi_{j}}}\\
&=\sum_{P}\zeta^{P}(H_{0}^{-1})_{P_{n},n}\dots(H_{0}^{-1})_{P_{1},1} \tag{5.4.1}
\end{align*}
其中 $\zeta=-1$。（若取 $\zeta=+1$，则上式便是玻色算符的 Wick 定理。）

让我们把上述形式体系应用于零温下的自由费米子系统。考虑如下路径积分：
\begin{align*}
Z_{0} &= \int \mathcal{D}\psi^{*}\mathcal{D}\psi\,\mathrm{e}^{\mathrm{i}\int_{0}^{\beta}\mathrm{d}t\,L_{0}}\\
L_{0} &= \sum_{\pmb{k}}\mathrm{i}\psi_{\pmb{k}}(t)^{*}\partial_{t}\psi_{\pmb{k}}(t)-\psi_{\pmb{k}}(t)^{*}\xi_{\pmb{k}}\psi_{\pmb{k}}(t)
\end{align*}
我们可以算出路径积分平均
\begin{equation*}
\mathrm{i}G_{\pmb{k}}(t,t')=\left\langle \psi_{\pmb{k}}(t)\psi_{\pmb{k}}^{*}(t')\right\rangle=\frac{1}{(-\mathrm{i})(\mathrm{i}\partial_{t}-\xi_{\pmb{k}})}
\end{equation*}
在频率空间中，我们有
\begin{equation*}
G_{\pmb{k},\omega}=\frac{1}{\omega-\xi_{\pmb{k}}}
\end{equation*}
它与式 (4.2.9) 中的编时平均\footnote{这里给出的计算是一种形式计算，未能重现重要的 $0^{+}$ 项。2.2.2 节讨论过如何产生 $0^{+}$ 项。}$\langle 0|T(\psi_{\pmb{k}}(t)\psi_{\pmb{k}}^{\dagger}(t'))|0\rangle$ 相一致。利用 Wick 定理，我们可以导出路径积分平均与编时平均之间如下更为普遍的关系：
\begin{equation*}
\left\langle \psi_{i_{1}}\dots\psi_{i_{n}}\psi_{j_{n}}^{*}\dots\psi_{j_{1}}^{*}\right\rangle=\frac{\mathrm{Tr}\left(T(\psi_{i_{1}}\dots\psi_{i_{n}}\psi_{j_{n}}^{\dagger}\dots\psi_{j_{1}}^{\dagger}\,\mathrm{e}^{-\int_{0}^{\beta}\mathrm{d}t\,H_{0}})\right)}{\mathrm{Tr}\left(T(\mathrm{e}^{-\int_{0}^{\beta}\mathrm{d}t\,H_{0}})\right)}
\end{equation*}
这样，我们就可以用路径积分计算任意的编时关联。

以上结果都是针对自由费米子系统的。让我们把路径积分应用于具有如下形式的相互作用系统：
\begin{equation*}
H=H_{0}+H_{I}(t)
\end{equation*}
其中 $H_{0}$ 描述自由费米子系统。设 $H_{I}(t)$ 是 $t$ 的缓变函数，在有限的 $t$ 处 $H_{I}(t)=H_{I}$，而 $H_{I}(\pm\infty)=0$。令
\begin{equation*}
Z=\left\langle 0\left|T\bigl(\mathrm{e}^{-\mathrm{i}\int \mathrm{d}t\,(H_{0}+H_{I}(t))}\bigr)\right|0\right\rangle
\end{equation*}
为配分函数，其中 $|0\rangle$ 是 $H_{0}$ 的基态。配分函数 $Z$ 有如下微扰展开：
\begin{align*}
Z &= Z_{0}\left(1+Z_{0}^{-1}\left\langle 0\left|T\left((- \mathrm{i}\int \mathrm{d}t\,H_{I})\,\mathrm{e}^{-\mathrm{i}\int \mathrm{d}t\,H_{0}}\right)\right|0\right\rangle+\dots\right)\\
Z_{0} &= \left\langle 0\left|T(\mathrm{e}^{-\mathrm{i}\int \mathrm{d}t\,H_{0}})\right|0\right\rangle.
\end{align*}
由于薛定谔绘景中的 $Z_{0}^{-1}\langle 0|T(O_{1}(t_{1})\dots\allowbreak O_{n}(t_{n})\,\mathrm{e}^{-\mathrm{i}\int \mathrm{d}t\,H_{0}})|0\rangle$ 等于海森堡绘景中的 $\langle 0|T(O_{1}(t_{1})\dots\allowbreak O_{n}(t_{n}))|0\rangle$，这两个量都表示编时平均。我们看到，相互作用系统的配分函数 $Z$ 可以用自由费米子系统 $H_{0}$ 中的编时平均表出。这使我们能够把 $Z$ 用路径积分表示如下：
\begin{equation*}
Z=\int \mathcal{D}\psi^{*}\mathcal{D}\psi\,\mathrm{e}^{\mathrm{i}\int_{0}^{\beta}\mathrm{d}t\,(L_{0}+L_{I})},\qquad L_{I}=-H_{I}
\end{equation*}
类似地，利用微扰展开可以证明，相互作用系统的（薛定谔绘景中的）编时关联可以用路径积分表示如下：
\begin{equation*}
\frac{\left\langle 0\left|T(O_{1}(t_{1})\dots O_{n}(t_{n})\,\mathrm{e}^{-\mathrm{i}\int \mathrm{d}t\,H})\right|0\right\rangle}{\left\langle 0\left|T(\mathrm{e}^{-\mathrm{i}\int \mathrm{d}t\,H})\right|0\right\rangle}=\frac{\int \mathcal{D}\psi^{*}\mathcal{D}\psi\,O_{1}(t_{1})\dots O_{n}(t_{n})\,\mathrm{e}^{\mathrm{i}\int_{0}^{\beta}\mathrm{d}t\,(L_{0}+L_{I})}}{\int \mathcal{D}\psi^{*}\mathcal{D}\psi\,\mathrm{e}^{\mathrm{i}\int_{0}^{\beta}\mathrm{d}t\,(L_{0}+L_{I})}}
\end{equation*}
现在，让我们用路径积分微扰地计算配分函数 $Z$：
\begin{align*}
Z &= Z_{0}\,\frac{\int \mathcal{D}\psi^{*}\mathcal{D}\psi\,\mathrm{e}^{\mathrm{i}\int \mathrm{d}t\,(L_{0}+L_{I})}}{\int \mathcal{D}\psi^{*}\mathcal{D}\psi\,\mathrm{e}^{\mathrm{i}\int \mathrm{d}t\,L_{0}}}\\
&=Z_{0}\left(1+\left\langle \left(\mathrm{i}\int \mathrm{d}t\,L_{I}\right)\right\rangle_{0}+\frac{1}{2!}\left\langle \left(\mathrm{i}\int \mathrm{d}t\,L_{I}\right)^{2}\right\rangle_{0}+\dots\right)
\end{align*}
若假设相互作用具有如下形式：
\begin{equation*}
\int \mathrm{d}t\,L_{I}=-\int \mathrm{d}x\,\mathrm{d}x'\,\psi^{*}(x^{\mu})\psi(x^{\mu})V(x^{\mu},x'^{\mu})\psi^{*}(x'^{\mu})\psi(x'^{\mu})
\end{equation*}

%FIG{5.7}: {对 $Z/Z_{0}$ 贡献 $V$ 一阶项的图。}

其中 $x^{\mu}=(t,x^{1},\dots,x^{d})$，$\mathrm{d}x=\mathrm{d}t\,\mathrm{d}^{d}\pmb{x}$，那么我们发现
\begin{align*}
\frac{Z}{Z_{0}} &= 1+\left\langle \int \mathrm{d}x\,\mathrm{d}x'\,\psi^{*}(x^{\mu})\psi(x^{\mu})[-\mathrm{i}V(x^{\mu},x'^{\mu})]\psi^{*}(x'^{\mu})\psi(x'^{\mu})\right\rangle_{0}+\dots\\
&=1+\int \mathrm{d}x\,\mathrm{d}x'\,(-)\mathrm{i}G_{0}(0)[-\mathrm{i}V(x^{\mu},x'^{\mu})](-)\mathrm{i}G_{0}(0)\\
&\quad +\int \mathrm{d}x\,\mathrm{d}x'\,(-)\mathrm{i}G_{0}(x^{\mu}-x'^{\mu})\,\mathrm{i}G_{0}(x'^{\mu}-x^{\mu})[-\mathrm{i}V(x^{\mu},x'^{\mu})]+\dots
\end{align*}
$V$ 的两个一阶项可以用图 5.7 所示的费曼图来表示。下面关于费米子的费曼规则与玻色子的几乎完全相同，唯一的区别是每个闭合圈额外贡献一个 $-1$ 因子。
\begin{enumerate}
\item[(a)] 费曼图中的每条线代表一个传播子 $G_{0}(x,x')$，其中 $x$ 与 $x'$ 是该线两端两个顶点的坐标。
\item[(b)] 每个顶点代表对该顶点位置的一个积分 $\int \mathrm{d}x$。
\item[(c)] 每条虚线代表相互作用势 $-\mathrm{i}V(x,x')$。
\item[(d)] 每个闭合费米圈贡献一个 $-1$ 因子。
\item[(e)] 连通图的值由规则 (a)--(d) 给出，非连通图的值由连通子图的乘积给出。
\end{enumerate}
费曼图抓住了微扰展开的结构。它们不仅给出 $V$ 量级的项，还给出 $V^{2}$ 量级的项以及所有其他更高阶的项。

在 $Z/Z_{0}$ 的微扰展开中，路径积分平均 $\langle\dots\rangle_{0}$ 同时包含连通图与非连通图。我们可以用连链展开把 $Z/Z_{0}$ 的展开改写为
\begin{equation*}
\ln\frac{Z}{Z_{0}}=\sum_{n=0}^{\infty}\frac{1}{n!}\left\langle \left(\mathrm{i}\int \mathrm{d}t\,L_{I}\right)^{n}\right\rangle_{0c}=\left\langle \mathrm{e}^{\mathrm{i}\int \mathrm{d}t\,L_{I}}\right\rangle_{0c} \tag{5.4.2}
\end{equation*}
其中路径积分平均 $\langle\dots\rangle_{0c}$ 只包含连通图。

\subsection*{问题 5.4.1}

\begin{enumerate}
\item[(a)] 计算 $Z/Z_{0}$ 微扰展开中 $V^{2}$ 量级的项，并画出相应的费曼图。
\item[(b)] 计算 $\ln(Z/Z_{0})$ 微扰展开中 $V^{2}$ 量级的项，并画出相应的费曼图。
\end{enumerate}

\subsection{自能与两体相互作用}

现在我们可以对如下相互作用电子系统微扰地计算电子格林函数了：
\begin{equation*}
H=\sum_{\pmb{k}}\xi_{\pmb{k}}c_{\alpha\pmb{k}}^{\dagger}c_{\alpha\pmb{k}}+\frac{1}{2\mathcal{V}}\sum_{\pmb{k},\pmb{k}',\pmb{q}}c_{\alpha\pmb{k}}^{\dagger}c_{\alpha,\pmb{k}-\pmb{q}}\,V_{\pmb{q}}\,c_{\beta\pmb{k}'}^{\dagger}c_{\beta,\pmb{k}'+\pmb{q}}
\end{equation*}
编时格林函数由如下路径积分平均给出：
\begin{equation*}
\mathrm{i}G(\pmb{x},t)\delta_{\alpha\beta}=\frac{\int \mathcal{D}c\,\mathcal{D}c^{*}\,c_{\alpha}(\pmb{x},t)c_{\beta}^{*}(0,0)\,\mathrm{e}^{\mathrm{i}\int \mathrm{d}t\,L}}{\int \mathcal{D}c\,\mathcal{D}c^{*}\,\mathrm{e}^{\mathrm{i}\int \mathrm{d}t\,L}}\equiv\left\langle c_{\alpha}(\pmb{x},t)c_{\beta}^{*}(0,0)\right\rangle
\end{equation*}
这是针对相互作用电子的，其中
\begin{equation*}
L=\mathrm{i}c_{\alpha}^{*}\partial_{t}c_{\alpha}-H
\end{equation*}
是相互作用电子的拉格朗日量。注意，这里我们假设了 $\langle c_{\alpha}(\pmb{x},t)c_{\beta}^{*}(0,0)\rangle$ 在自旋转动下不变，因而它具有 $\mathrm{i}G(\pmb{x},t)\delta_{\alpha\beta}$ 的形式。

为计算 $\mathrm{i}G$，方便的做法是引入
\begin{equation*}
Z[\eta_{\alpha},\eta_{\alpha}^{*}]=\int \mathcal{D}c\,\mathcal{D}c^{*}\,\mathrm{e}^{\mathrm{i}\int \mathrm{d}t\,L+\int \mathrm{d}t\,\mathrm{d}^{d}\pmb{x}\,(\eta_{\alpha}c_{\alpha}^{*}+\eta_{\alpha}^{*}c_{\alpha})}
\end{equation*}
其中 $\eta(\pmb{x},t)$ 与 $\eta^{*}(\pmb{x},t)$ 同 $c(\pmb{x},t)$ 与 $c^{*}(\pmb{x},t)$ 一样，都是格拉斯曼数。我们看到
\begin{equation*}
\mathrm{i}G(\pmb{x},t)\delta_{\alpha\beta}=\frac{\partial}{\partial\eta_{\alpha}^{*}(\pmb{x},t)}\frac{\partial}{\partial\eta_{\beta}(0,0)}\ln Z[\eta_{\alpha},\eta_{\alpha}^{*}]
\end{equation*}
利用式 (5.4.2)，我们得到
\begin{align*}
\mathrm{i}G(\pmb{x},t)\delta_{\alpha\beta} &= \frac{\partial}{\partial\eta_{\alpha}^{*}(\pmb{x},t)}\frac{\partial}{\partial\eta_{\beta}(0,0)}\left\langle \mathrm{e}^{\mathrm{i}\int \mathrm{d}t\,L_{I}+\int \mathrm{d}t\,\mathrm{d}^{d}\pmb{x}\,(\eta_{\alpha}c_{\alpha}^{*}+\eta_{\alpha}^{*}c_{\alpha})}\right\rangle_{0c}\\
&=\left\langle c_{\alpha}(\pmb{x},t)c_{\beta}^{*}(0,0)\,\mathrm{e}^{\mathrm{i}\int \mathrm{d}t\,L_{I}}\right\rangle_{0c} \tag{5.4.3}
\end{align*}
这个公式为用费曼图计算相互作用电子的编时格林函数奠定了基础。

%FIG{5.8}: {对 $\mathrm{i}G_{\pmb{k}\omega}$ 有贡献的两个费曼图。虚线代表势 $-\mathrm{i}V_{\pmb{q}}$。}

准确到 $V_{\pmb{q}}$ 的一阶，我们得到
\begin{equation*}
\mathrm{i}G(\pmb{x},t)\delta_{\alpha\beta}=\left\langle c_{\alpha}(\pmb{x},t)c_{\beta}^{\dagger}(0,0)\right\rangle_{0}+\left\langle c_{\alpha}(\pmb{x},t)c_{\beta}^{\dagger}(0,0)(-\mathrm{i})\int \mathrm{d}t'\,H_{I}(t')\right\rangle_{0c}
\end{equation*}
利用 Wick 定理，我们可以把上式用自由电子格林函数 $G_{0}$ 表出。在 $\pmb{k}$--$\omega$ 空间中，我们得到（见问题 5.4.2）
\begin{align*}
\mathrm{i}G_{\pmb{k}\omega} &= \mathrm{i}G_{0,\pmb{k}\omega}+(-1)(\mathrm{i}G_{0,\pmb{k}\omega})^{2}(-\mathrm{i}V_{0})\frac{1}{\mathcal{V}}\sum_{\pmb{k}',\alpha'}\int \frac{\mathrm{d}\nu}{2\pi}\,\mathrm{i}G_{0,\pmb{k}'\nu}\,\mathrm{e}^{\mathrm{i}0^{+}\nu}\\
&\quad +(\mathrm{i}G_{0,\pmb{k}\omega})^{2}\frac{1}{\mathcal{V}}\sum_{\pmb{q}}\int \frac{\mathrm{d}\nu}{2\pi}\,\mathrm{i}G_{0,\pmb{k}-\pmb{q},\nu}\,(-\mathrm{i}V_{\pmb{q}})\,\mathrm{e}^{-\mathrm{i}0^{+}\nu} \tag{5.4.4}
\end{align*}
其中自由电子格林函数 $G_{0,\pmb{k}\omega}$ 由式 (4.2.4) 给出。两个 $V_{\pmb{q}}$ 量级的项总结在图 5.8 所示的费曼图中。动量空间中的费曼规则与实空间中的不同，如下所示。
\begin{enumerate}
\item[(a)] 费曼图中的每条线携带一个能量--动量矢量，比如 $(\pmb{q},\nu)$。它代表一个传播子 $\mathrm{i}G_{\pmb{q}\nu}=\mathrm{i}/(\nu-\xi_{\pmb{q}})$。
\item[(b)] 流入一个节点的能量--动量是守恒的。
\item[(c)] 每个圈代表积分 $(\beta\mathcal{V})^{-1}\int \frac{\mathrm{d}\nu}{2\pi}\sum_{\pmb{q}}$（译注：原书误排为 $\mathrm{d}t$）。
\item[(d)] 虚线代表相互作用 $-\mathrm{i}V_{\pmb{q}}$，其中 $\pmb{q}$ 是流过虚线的动量。
\item[(e)] 由于费米子算符的反对易性质，每个闭合圈额外贡献一个 $-1$ 因子。（玻色系统的动量空间费曼规则没有这样的 $-1$ 因子。）
\end{enumerate}

式 (5.4.4) 可以写成
\begin{equation*}
\mathrm{i}G_{\pmb{k}\omega}=\mathrm{i}G_{0,\pmb{k}\omega}+(\mathrm{i}G_{0,\pmb{k}\omega})^{2}(-\mathrm{i}\Sigma_{\pmb{k}\omega})
\end{equation*}

%FIG{5.9}: {对自能 $-\mathrm{i}\Sigma_{\pmb{k}\omega}$ 有贡献的两个图。}

%FIG{5.10}: {对 $\mathrm{i}G_{\pmb{k}\omega}$ 有贡献的更高阶图。}

其中
\begin{align*}
&-\mathrm{i}\Sigma_{\pmb{k}\omega}\\
&=(-1)\frac{-\mathrm{i}V_{0}}{\mathcal{V}}\sum_{\pmb{k}',\alpha'}\int \frac{\mathrm{d}\nu}{2\pi}\,\mathrm{i}G_{0,\pmb{k}'\nu}\,\mathrm{e}^{\mathrm{i}0^{+}\nu}+\frac{1}{\mathcal{V}}\sum_{\pmb{q}}\int \frac{\mathrm{d}\nu}{2\pi}\,\mathrm{i}G_{0,\pmb{k}-\pmb{q},\nu}\,(-\mathrm{i}V_{\pmb{q}})\,\mathrm{e}^{-\mathrm{i}0^{+}\nu}.
\end{align*}
上式可以用图 5.9 中的图来表示。如果我们把图 5.10 中的图所代表的某些更高阶项包括进来，格林函数就可以写成
\begin{equation*}
\mathrm{i}G_{\pmb{k}\omega}=\mathrm{i}G_{0,\pmb{k}\omega}\sum_{n=0}^{\infty}\left[(\mathrm{i}G_{0,\pmb{k}\omega})(-\mathrm{i}\Sigma_{\pmb{k}\omega})\right]^{n}
\end{equation*}
我们得到
\begin{equation*}
G_{\pmb{k}\omega}=\frac{1}{\omega-\xi_{\pmb{k}}-\Sigma_{\pmb{k}\omega}+\mathrm{i}\,\mathrm{sgn}(\omega)0^{+}}
\end{equation*}
色散关系既可以由 $G_{\pmb{k}\omega}$ 的极点位置得到，也可以由 $\omega-\xi_{\pmb{k}}-\Sigma_{\pmb{k}\omega}$ 的零点位置得到。我们看到，$\Sigma_{\pmb{k}\omega}$ 修正了电子（或准粒子）的能量，被称为自能。

注意
\begin{equation*}
\int \frac{\mathrm{d}\nu}{2\pi}\,\mathrm{i}G_{0,\pmb{k}'\nu}\,\mathrm{e}^{\mathrm{i}0^{+}\nu}=\mathrm{i}\int \frac{\mathrm{d}\nu}{2\pi}\,\frac{1}{\nu-\xi_{\pmb{k}'}+\mathrm{i}0^{+}\mathrm{sgn}(\nu)}\,\mathrm{e}^{\mathrm{i}0^{+}\nu}=-\Theta(-\xi_{\pmb{k}'})=-n_{\pmb{k}'}
\end{equation*}
由于围道必须在复 $\nu$ 平面的上半平面闭合，只有 $\nu$ 为负的极点才有贡献。由于 $\sum_{\pmb{k}',\alpha}\Theta(-\xi_{\pmb{k}'})$ 恰好就是电子的总数，自能中的第一项正是 Hartree 贡献（图 5.9(a)）：
\begin{equation*}
\Sigma_{\pmb{k}\omega}=V_{0}\rho_{0}+\dots
\end{equation*}
类似地，自能中的第二项（图 5.9(b)）具有如下形式：
\begin{equation*}
\Sigma_{\pmb{k}\omega}=\dots+\frac{1}{\mathcal{V}}\sum_{\pmb{q}}(-n_{\pmb{k}-\pmb{q}})V_{\pmb{q}}
\end{equation*}
这正是自能中的 Fock 贡献。

%FIG{5.11}: {图 (a)--(e) 对 $V_{\pmb{q}}^{2}$ 量级的自能有贡献。图 (f) 对自能没有贡献，因为它已包含在图 5.10 之中。}

在上面，我们看到了如何从这些图计算自能 $\Sigma_{\pmb{k}\omega}$ 以及准粒子色散 $\xi_{\pmb{k}}^{*}=\xi_{\pmb{k}}+\Sigma_{\pmb{k}\omega}\big|_{\omega=\xi_{\pmb{k}}^{*}}$。这一做法在 $V$ 的一阶上重现了 Hartree--Fock 结果。而图方法还使我们能够以系统的方式计算自能的更高阶贡献（见图 5.11）。

我们还可以用图来计算费米液体理论中的费米液体函数。我们知道，相互作用 $V_{\pmb{q}}$ 可以把动量为 $\pmb{k}_{1}$ 和 $\pmb{k}_{2}$ 的两个电子散射为动量为 $\pmb{k}_{1}'$ 和 $\pmb{k}_{2}'$ 的两个电子。费米液体理论中的相互作用项 $n_{\pmb{k}_{1}}n_{\pmb{k}_{2}}V(\pmb{k}_{1}-\pmb{k}_{2})$ 对应于把动量为 $\pmb{k}_{1}$ 和 $\pmb{k}_{2}$ 的两个电子散射为具有相同动量的两个电子的那些项。于是，费米液体函数 $f$ 接收两个贡献，由图 5.12 中的两个图所示：
\begin{equation*}
-\mathrm{i}f(\pmb{k}_{1},\alpha;\pmb{k}_{2},\beta)=(-\mathrm{i}V_{0})+(-)(-\mathrm{i}V_{\pmb{q}})\delta_{\alpha\beta}
\end{equation*}

%FIG{5.12}: {对费米液体函数 $f(\pmb{k}_{1},\alpha;\pmb{k}_{2},\beta)$ 有贡献的图。(a) 这里 $\pmb{q}=\pmb{k}_{1}-\pmb{k}_{1}=0$ 是动量转移，$\nu=\omega_{1}-\omega_{1}'$ 是能量转移。(b) 这里 $\pmb{q}=\pmb{k}_{1}-\pmb{k}_{2}$ 是动量转移，$\nu=\omega_{1}-\omega_{2}'$ 是能量转移。}

$(-)$ 号来自交换，而 $\delta_{\alpha\beta}$ 来自散射前后两个电子必须具有相同自旋这一要求。由于这里考虑的相互作用是瞬时的，$V_{\pmb{q}}$ 不依赖于能量转移 $\nu$（见图 5.12）。由此得到的费米液体函数 $f$ 也是瞬时的，且不依赖于能量转移 $\nu$：
\begin{equation*}
f(\pmb{k}_{1},\alpha;\pmb{k}_{2},\beta)=V_{0}-V_{\pmb{k}_{1}-\pmb{k}_{2}}\delta_{\alpha\beta}
\end{equation*}
这正是我们之前得到的 Hartree--Fock 结果 (5.3.2)。

\subsection*{问题 5.4.2}

用微扰论证明式 (5.4.4)。我们注意到，$\int \frac{\mathrm{d}\nu}{2\pi}\,\allowbreak\mathrm{i}G_{0,\pmb{k}'\nu}\,\allowbreak\mathrm{e}^{\mathrm{i}0^{+}\nu}$ 代表等时下的 $\langle c^{\dagger}c\rangle$，而 $\int \frac{\mathrm{d}\nu}{2\pi}\,\allowbreak\mathrm{i}G_{0,\pmb{k}'\nu}\,\allowbreak\mathrm{e}^{-\mathrm{i}0^{+}\nu}$ 代表等时下的 $\langle cc^{\dagger}\rangle$。

\subsection*{问题 5.4.3}

求自能 $\Sigma_{\pmb{k}\omega}$ 的微扰展开中 $V_{\pmb{q}}^{2}$ 量级各项的表达式，即图 5.11(a)--(e) 所描述的那些项。

\subsection{无规相位近似与有效势}

\begin{keypoints}
\item 粒子--空穴激发对相互作用势的屏蔽可以包含在无规相位近似（random phase approximation, RPA）之中。
\end{keypoints}

我们已经看到，在三维情形且相互作用为库仑相互作用时，Hartree--Fock 近似给出的自能在费米面附近是奇异的。这一奇异性导致无穷大的费米速度。奇异性来自库仑相互作用的长程势。短程相互作用不能产生任何奇异的自能。然而，我们知道，金属中的
